{\Huge \bfseries \noindent Chapter 6}
\vspace{-1em} 
\section*{\LARGE Discussion}
\addcontentsline{toc}{section}{6\quad Discussion}
\setcounter{section}{6}
\setcounter{subsection}{0}
\renewcommand{\chapternum}{6}

% Reset figure/table numbering per chapter
\renewcommand{\thefigure}{\chapternum.\arabic{figure}}
\setcounter{figure}{0}
\renewcommand{\thetable}{\chapternum.\arabic{table}}
\setcounter{table}{0}

\vspace{2em}

\subsection{Overview}
This thesis set out to address two interlinked obstacles that have constrained the study of circRNAs: the computational difficulty of detecting, quantifying, and reliably annotating circRNAs from high-throughput sequencing data, and the incomplete understanding of their post-transcriptional regulatory functions in oncogenesis and drug resistance. To this end, the work presented herein delivers two open-source software contributions to the bioinformatics and RNA biology communities alongside a translational investigation into the regulatory roles of circRNAs in PCa enzalutamide resistance and BCR, through the lens of the ceRNA network. circRNAs represent compelling candidates for both biomarker discovery and therapeutic intervention owing to their covalently closed loop structure, which lacks accessible 5' and 3' ends, conferring intrinsic resistance to exonucleolytic degradation and a markedly extended half-life relative to linear transcripts \cite{Salzman2012, Li2015Aug}.

In Chapter One, the biological foundation of androgen receptor signalling in prostate development, tissue homeostasis, and malignant transformation was described \cite{Rebello2021Feb}. While androgen-deprivation therapy and second-generation anti-androgens such as enzalutamide initially produce clinical responses, the selective pressure of AR blockade inevitably precipitates treatment failure, driving progression to castration-resistant prostate cancer and subsequent biochemical recurrence \cite{Scher2012Sep, Cornford2017, Rans2022, Shore2024}. The molecular mechanisms driving this progression—including AR gene amplification, point mutations within the ligand-binding domain, and the emergence of constitutively active splice variants such as \textit{AR-V7} \cite{Azad2015May, Hu2009Jan, Csizmarik2021Oct}—underscore the necessity for prognostic biomarkers capable of capturing latent therapeutic resistance.

Circular RNAs frequently harbour evolutionarily conserved microRNA response elements within their mature spliced sequences, enabling them to competitively bind and sequester miRNAs and relieve post-transcriptional repression on downstream messenger RNA targets within competing endogenous RNA networks \cite{Hansen2013Mar, Thomas2014Aug, Denzler2014Jun}. This regulatory interplay makes circRNAs attractive candidates for untangling complex oncogenic phenotypes such as therapeutic resistance. However, translating this mechanistic premise into reproducible prognostic models has historically been obstructed by bioinformatics bottlenecks: high discordance among back-splice junction detection algorithms, the absence of standardised quantification workflows, and fragmented public databases with conflicting coordinate systems \cite{Chen2023Jan, Digby2023Dec, Digby2024Dec}. This thesis systematically addresses these challenges across four research chapters. By first diagnosing tool-level discrepancies (Chapter Two), developing an automated consensus-based quantification pipeline in nf-core/circrna (Chapter Three), and creating a unified multi-database annotation resource in pycircdb (Chapter Four), these software contributions establish the scalable infrastructure needed to support complex circRNA analyses, complementing the translational investigation in Chapter Five where an enzalutamide-resistance ceRNA network was modelled to derive a 23-gene prognostic signature for biochemical recurrence in prostate cancer.

\subsection{Synopsis of Major Findings}
This thesis evaluated the computational landscape of circRNA detection, delivered two novel software frameworks that resolve challenges in circRNA quantification and annotation, and investigated the regulatory mechanism of non-coding RNAs in prostate cancer drug resistance. The principal findings of this work are:

\begin{itemize}
    \item \textbf{The computational landscape of circRNA detection is fragmented and lacks a gold standard.} A comprehensive evaluation of the computational strategies underpinning circRNA analysis - ranging from library enrichment protocols (such as ribonuclease R digestion and RPAD) to segmented-based and pseudo-reference mapping algorithms - revealed substantial discordance among established detection tools \cite{Panda2017Jul, Pandey2019Feb, Hansen2016Apr, Zeng2017Jun, Vromman2023Aug}. Benchmarking evidence demonstrated that individual tools exhibit high tool-specific false-positive rates and divergent sensitivities depending on aligner choices and filtering parameters. The findings strongly indicated that multi-tool consensus calling represents the most effective strategy to mitigate algorithmic blind spots, while identifying the absence of standardised coordinate conventions and nomenclature as a critical barrier to cross-study reproducibility \cite{CirComPara2, Chen2023Jan}.

    \item \textbf{nf-core/circrna delivers an automated, containerised, consensus-driven pipeline for end-to-end circRNA analysis.} To address the workflow fragmentation and tool discordance identified in Chapter Two, this work introduced nf-core/circrna, the first workflow to integrate circRNA quantification across seven tools, automated miRNA target prediction, and differential expression analysis into a single Nextflow DSL2 pipeline \cite{Digby2023Dec, DiTommaso2017, Ewels2020Mar}. Combinatorial evaluation of nf-core/circrna's seven quantification tools demonstrated that an ensemble of three tools requiring consensus support from at least two tools achieves an optimal trade-off between sensitivity ($97.2\%$), precision ($98.5\%$), and computational efficiency ($F_1 = 97.8\%$). Furthermore, benchmarking highlighted the robust individual performance of CIRIquant \cite{CIRIquant} and identified an elevated false-positive rate associated with DCC \cite{DCC} due to unconstrained STAR 2-pass chimeric mapping.

    \item \textbf{pycircdb unifies fragmented circRNA databases for high-throughput programmatic annotation.} To resolve the nomenclature and annotation bottlenecks identified throughout the literature \cite{Chen2023Jan}, this thesis delivered pycircdb, a Python-based command-line tool that harmonises nine public circRNA databases into a single queryable system \cite{Digby2024Dec}. pycircdb unifies genomic coordinates to a 0-based format, resolves hg19-hg38 liftovers while preserving segmental duplications, and uses Apache Hamilton with columnar Parquet storage. Consequently, it enables rapid batch annotation of candidate circRNAs with precomputed miRNA and RBP interactions on modest hardware (processing 300,000 circRNAs in 66.7s with 3.8 GB peak RAM and 4.2 GB disk footprint including database download).

    \item \textbf{A competing endogenous RNA network-derived signature predicts five-year biochemical recurrence in prostate cancer.} By integrating transcriptomic data from enzalutamide-resistant versus enzalutamide-sensitive LNCaP cell lines with clinical PCa patient cohorts across circRNA, miRNA, and mRNA biotypes, this work constructed a biologically constrained ceRNA network. This knowledge-based feature selection framework yielded a 23-gene prognostic signature predictive of five-year biochemical recurrence. The signature demonstrated robust discriminative performance across nine independent cohorts ($n=1{,}451$), achieving an apparent Harrell's C-index of 0.73, an optimism-corrected C-index of 0.68 via bootstrap resampling ($M=200$), and strong external validation ($C = 0.71$, five-year time-dependent $\text{AUC} = 0.87$) \cite{Mortensen2015Nov, Harrell1996Feb, Steyerberg2015Apr}. Crucially, the 23-gene signature was found to comprise a largely unique gene set relative to 30 published prognostic signatures, underscoring the discovery value of regulatory network modelling.
\end{itemize}

\subsection{Synthesis of Research Themes}

\subsubsection{Scalable, reproducible, and standardised workflows for circRNA analyses}
The computational challenges inherent in characterising circRNAs reflect a broader reproducibility crisis across high-throughput genomics and transcriptomics studies. Historically, computational pipelines frequently relied on bespoke, unversioned in-house scripts, non-containerised execution environments, and manual file formatting, leading to software decay and severe barriers to experimental replication \cite{Ziemann2023Oct}. In circRNA analyses, this fragmentation was compounded by the algorithmic divergence between segmented-anchor approaches (such as find\_circ and UROBORUS) and pseudo-reference methods (such as KNIFE and CIRIquant), where aligner-specific indexing heuristics, anchor-extension rules, and back-splice junction filtering thresholds produce widely divergent candidate sets across tools \cite{findcirc, UROBORUS, KNIFE, CIRIquant, Hansen2016Apr, Zeng2017Jun, Vromman2023Aug}. Furthermore, downstream public circRNA databases emerged as isolated web repositories with conflicting coordinate systems (0-based versus 1-based), asynchronous genome builds (hg19 and hg38), and idiosyncratic identifier schemes, hindering cross-database interoperability \cite{Chen2023Jan}.

In response to widespread reproducibility challenges in computational science, modern bioinformatics has undergone a transformative paradigm shift toward declarative, domain-specific workflow management systems, most notably Nextflow \cite{DiTommaso2017} and Snakemake \cite{Koster2012Oct, Molder2021Feb}. These systems decouple pipeline logic from compute infrastructure, allowing automated dependency resolution and scalable execution across local machines, high-performance compute clusters (such as SLURM, SGE or Kubernetes), and commercial cloud environments. In parallel, software containerisation through Docker, Apptainer (formerly Singularity), and community repositories such as BioContainers warehouse pinned versions of bioinformatics tools as container images \cite{Docker, Apptainer, daVeigaLeprevost2017Aug}, enabling reproducible analyses and alleviating dependency conflicts that have historically hindered analysis. These technical advances have been formalised through international standards and open-source initiatives, including the Global Alliance for Genomics and Health (GA4GH) Tool Registry and Workflow Execution Services \cite{Rehm2021Nov}, the Common Workflow Language (CWL) \cite{Crusoe2022Jun}, and centralised workflow registries such as WorkflowHub \cite{Goble2021Oct}, all driving adherence to the FAIR (Findable, Accessible, Interoperable, Reusable) data stewardship principles \cite{Wilkinson2016Mar}.

The computational contributions delivered in this thesis—\texttt{nf-core/circrna} (Chapter Three) and \texttt{pycircdb} (Chapter Four)—implement these modern workflow and data standardisation principles to establish an integrated analytical foundation for circRNA biology \cite{Digby2023Dec, Digby2024Dec}. Rather than proposing yet another detection algorithm, \texttt{nf-core/circrna} was engineered within the nf-core framework \cite{Ewels2020Mar}, embedding continuous integration testing, pinned containerised dependencies, and standardised configuration profiles across seven distinct circRNA quantification tools. By automating raw read quality control, multi-tool alignment, consensus-based back-splice junction filtering, mature spliced sequence extraction, miRNA response element prediction, and differential expression analysis into a single, platform-agnostic pipeline, \texttt{nf-core/circrna} mitigates the algorithmic blind spots and high false-positive rates inherent to single-tool analyses while ensuring deterministic execution across diverse computing platforms.

Complementing this upstream quantification pipeline, \texttt{pycircdb} resolves the persistent database fragmentation and annotation bottlenecks in the circRNA field. Moving away from relational database installations and static web interfaces that cannot scale to high-throughput batch queries, \texttt{pycircdb} restructures nine public repositories into a unified, 0-based coordinate framework. By leveraging columnar Apache Parquet storage hosted on public cloud infrastructure alongside Apache Hamilton directed acyclic graphs for parallel execution, \texttt{pycircdb} enables rapid, partition-aware batch queries of circRNA genomic annotations, spliced FASTA sequences, and precomputed miRNA and RBP interactions directly from the command line on standard workstation hardware. Together, \texttt{nf-core/circrna} and \texttt{pycircdb} advance computational reproducibility and lower technical barriers, bridging the gap between raw sequencing reads, circRNA quantification and standardised functional annotation.

\subsubsection{circRNAs as regulatory agents in the competing endogenous RNA network}
Complex human pathologies, particularly cancers, do not arise from isolated genetic abnormalities acting in a linear hierarchy; rather, they reflect systemic perturbations across interconnected networks \cite{Barabasi2011Jan, Ideker2012Jan}. Network biology posits that cellular phenotypes are governed by multi-layered regulatory graphs spanning protein-protein interactions, transcriptional circuits, and post-transcriptional cascades \cite{Califano2011Jan, Anastasiadou2018Jan}. Within this systems biology framework, the competing endogenous RNA (ceRNA) hypothesis formulated by Salmena et al. \cite{Salmena2011Aug} and expanded by Tay et al. \cite{Tay2014Aug} revealed a hidden post-transcriptional regulatory language: RNA transcripts harbouring shared microRNA response elements cross-regulate one another by competitively binding the limited intracellular pool of mature miRNAs. The subsequent discovery of circRNAs added an exceptionally stable, potent regulatory layer to this network architecture. Owing to their covalently closed loop structure, circRNAs resist exonucleolytic degradation by RNase R, conferring extended intracellular half-lives and enabling them to act as persistent molecular sponges that sequester miRNAs away from protein-coding mRNAs \cite{Salzman2012, Li2015Aug, Hansen2013Mar, Denzler2014Jun, Conn2024Sep}. This titrating activity coordinates critical oncogenic programmes, including cell proliferation, epithelial-to-mesenchymal transition, and acquired therapeutic resistance across diverse human malignancies \cite{Kristensen2019Nov, Hsiao2017, Lim2021Dec}.

In translational oncology, high-dimensional transcriptomic profiling regularly confronts the curse of dimensionality, where the number of measured genomic features vastly outweighs the number of available clinical patient samples. When using unconstrained data-driven feature selection or standard machine learning algorithms, models are acutely prone to overfitting technical noise, library batch effects, and platform-specific artefacts inherent to retrospective multicentre patient cohorts \cite{Riley2016Jun, Dhiman2021Oct, Steyerberg2015Apr}. In Chapter Five, this limitation was addressed by employing the topological constraints of the multi-layered ceRNA network as an explicit, knowledge-based inductive prior. Rather than selecting candidate prognostic biomarkers solely through unconstrained differential expression or univariate screening, candidate features were filtered to enforce coherent, reciprocal expression trajectories across circRNA, miRNA, and mRNA biotypes (upregulated circRNA $\rightarrow$ downregulated miRNA $\rightarrow$ upregulated mRNA, and vice versa).

Embedding this multi-omic regulatory network structure into the feature selection pipeline effectively regularised the high-dimensional search space prior to penalised survival modelling. As shown by the resulting 23-gene prognostic signature for prostate cancer biochemical recurrence, grounding biomarker discovery via a biologically constrained network topology insulated the model against cohort-specific batch variance. This inductive network resulted in a clinically transportable prognostic signal that generalised across diverse international patient cohorts evaluated under internal-external cross-validation, bootstrap resampling, and independent external validation \cite{Mortensen2015Nov, Harrell1996Feb, Steyerberg2015Apr}. Strikingly, the network-derived gene set demonstrated minimal overlap with thirty published prostate cancer prognostic panels, capturing a largely distinct transcriptional subspace \cite{Knezevic2013Oct, Cuzick2011Mar, Long2014Jun}. This highlights the discovery power of systems-level network modelling: by leveraging non-coding regulatory interactions as functional filters, researchers can uncover stable, clinically actionable biomarker programmes that conventional gene-centric and data-driven screening pipelines might overlook.

\subsubsection{Androgen receptor signalling plasticity, drug resistance and biomarker discovery}
Prostate cancer progression is fundamentally driven by the androgen receptor (AR) signalling axis. As detailed in Chapter One, canonical AR signalling governs prostate development, epithelial differentiation, and secretory function \cite{Rebello2021Feb}. Consequently, androgen-deprivation therapy (ADT) and surgical ablation have long served as the primary therapeutic action for localised and advanced disease. However, under the selective pressure of systemic androgen blockade, tumour cells undergo evolutionary adaptation, progressing to castration-resistant prostate cancer (CRPC) \cite{Scher2012Sep, Cornford2017}. Although second-generation anti-androgens such as enzalutamide competitively bind the AR ligand-binding domain to extend patient survival, intrinsic resistance occurs in up to 40\% of patients, and acquired resistance inevitably develops, heralding biochemical recurrence and metastatic dissemination \cite{Rans2022, Shore2024}.

The molecular underpinnings of enzalutamide resistance reflect transcriptional and genomic plasticity. Resistant tumour clones circumvent therapeutic inhibition through several distinct mechanisms, including AR gene amplification, point mutations within the ligand-binding domain (such as T877A, H874Y, and F876L, which mediate antagonist-to-agonist switching), and the generation of constitutively active splice variants lacking the ligand-binding domain, such as \textit{AR-V7} \cite{Azad2015May, Hu2009Jan, Csizmarik2021Oct}. Chapter Five leveraged in vitro cell-line models of acquired enzalutamide resistance (LNCaP clone models) to investigate transcriptomic alterations that accompany sustained anti-androgen exposure, capturing non-coding and coding RNA perturbations associated with treatment failure.

By integrating transcriptomic profiles from these cell-line models with PCa patient survival cohorts, Chapter Five assessed whether molecular programmes associated with acquired enzalutamide resistance could inform postoperative clinical risk stratification. Evaluating the 23-gene panel across retrospective cohorts demonstrated that transcriptional changes induced during acquired enzalutamide resistance are statistically informative for predicting time to biochemical recurrence following primary definitive therapy. Furthermore, several individual genes within the signature represent functionally relevant candidates in therapy resistance and tumour progression, including \textit{SFRP4}, \textit{B4GALT1}, \textit{FA2H}, and \textit{NEU3} \cite{Sandsmark2017Oct, DeVitis2019Nov, Yao2019Mar, Kawamura2012Jan}. While prospective clinical validation remains necessary, this work illustrates the feasibility of using in vitro drug-resistance signals, constrained by regulatory RNA networks, to guide the prioritisation of candidate prognostic biomarkers for postoperative disease monitoring.

\subsection{Methodological and Biological Limitations}

\subsubsection{Computational challenges in circRNA isoform resolution and cross-database mapping}
While the computational methodologies developed in this thesis advance circRNA detection and annotation, several limitations must be acknowledged:

\begin{itemize}
    \item \textbf{Short-read back-splice junction limitations}: Both nf-core/circrna and the primary databases integrated in pycircdb rely primarily on short-read Illumina sequencing data. Short reads span the back-splice junction (BSJ) but cannot resolve the internal exon composition or full-length mature sequence of circRNAs that contain multiple alternatively spliced internal exons or retained introns (EIcircRNAs) \cite{CIRI-full, isoCirc}. Consequently, predicted miRNA response elements derived from reference-guided exon concatenation may overestimate target sites if specific internal exons are skipped in vivo.
    
    \item \textbf{Upstream database dependency and longevity}: pycircdb aggregates nine curated resources, but the tool remains inherently dependent on the accuracy, curation quality, and ongoing availability of the underlying source files. Several circRNA databases originally surveyed in Chapter Two (e.g., CircRiC, Circ2Disease, Circ2Traits) have become unmaintained or inaccessible, reflecting a broader challenge of resource sustainability in bioinformatics. Furthermore, pycircdb is currently constrained to human genome builds (hg19 and hg38), limiting its utility for model organisms.

    \item \textbf{Coordinate liftover asymmetries}: Although pycircdb explicitly handles one-to-many mappings arising from segmental duplications and repetitive loci during UCSC liftOver between hg19 and hg38, a small fraction of circRNA coordinates situated in complex genomic regions fail cross-assembly liftover, resulting in build-specific records.
\end{itemize}

\subsubsection{Model limitations and cohort heterogeneity in prostate cancer prognosis}
The prognostic modelling framework established in Chapter Five provides valuable biological candidates, but exhibits important constraints:

\begin{itemize}
    \item \textbf{Clinical metadata constraints across historical cohorts}: In the meta-analysis of ten retrospective prostate cancer cohorts, comprehensive clinical metadata—such as pathological tumour stage, surgical margin status, and lymph node involvement—were inconsistently recorded across public repositories. Consequently, multivariate adjustments during feature selection and model evaluation were limited to serum PSA and Gleason score, the two parameters universally available across all derivation studies.
    
    \item \textbf{Between-study heterogeneity and calibration}: Internal-external cross-validation revealed moderate between-study heterogeneity in calibration slopes ($I^2 = 43.3\%$), driven by variations in microarray/sequencing platforms and patient demographic differences across international cohorts. Attempts to directly model between-study variation using multivariable mixed-effects models with random slopes proved computationally intractable due to convergence failure, reflecting the substantial sample sizes and BCR events required to reliably estimate high-dimensional random-effects covariance structures. 

    \item \textbf{In vitro modelling of drug resistance}: The ceRNA network was in part derived using enzalutamide-resistant LNCaP cell-line models in 2D monoculture. While these models effectively isolate cell-autonomous, AR-dependent resistance mechanisms, they cannot capture the complex stromal-epithelial interactions, immune infiltration, hypoxia, or lineage plasticity (such as neuroendocrine transdifferentiation) present in primary tumour microenvironments.

    \item \textbf{Requirement for functional wet-lab validation}: Although the 23-gene signature identified robust candidates with strong statistical support, the precise causal mechanisms by which these genes modulate biochemical recurrence and anti-androgen resistance require empirical validation through targeted CRISPR-Cas knockout, overexpression, and circRNA sponge rescue experiments.
\end{itemize}

\subsection{Future Perspectives}

\subsubsection{Next-generation long-read and spatial circRNA profiling}
The rapid evolution of sequencing technologies presents exciting avenues for expanding the computational tools developed in this thesis:

\begin{itemize}
    \item \textbf{Long-read full-length circRNA workflows}: Third-generation long-read sequencing technologies, including Oxford Nanopore Technologies (ONT) and Pacific Biosciences (PacBio), enable end-to-end sequencing of intact circular transcripts following rolling-circle amplification (RCA) or reverse transcription (RCRT) \cite{isoCirc, CIRIlong, circFLseq}. Future development of the nf-core ecosystem should incorporate a dedicated long-read circRNA workflow to reconstruct complete circular transcriptomes, resolving internal alternative splicing events and microRNA binding site architecture without reference-guided imputation.
    
    \item \textbf{Expansion of pycircdb across model organisms and single-cell modalities}: Future iterations of pycircdb will extend coordinate reconciliation and lookup capabilities to major model organisms (including \textit{Mus musculus}, \textit{Rattus norvegicus}, and \textit{Drosophila melanogaster}), matching the multi-species capacity of nf-core/circrna. Additionally, incorporating emerging single-cell and spatial transcriptomic circRNA data will allow researchers to query cell-type-specific and spatially resolved circRNA expression across complex tissues.

    \item \textbf{Integration with deep learning splicing predictors}: Integrating deep-learning models of non-canonical splicing and RNA secondary structure (e.g., JEDI, CIRI-deep) into nf-core/circrna will improve de novo back-splice junction classification, particularly in low-coverage or poly(A)-selected libraries \cite{JEDI, CIRIdeep}.
\end{itemize}

\subsubsection{Clinical translation, liquid biopsies and functional characterisation}
From a translational perspective, the findings presented in this thesis open several promising directions for clinical oncology:

\begin{itemize}
    \item \textbf{Prospective validation of the 23-gene prognostic signature}: The 23-gene signature derived in Chapter Five warrants further evaluation across emerging, contemporary clinical study datasets. Assessing its prognostic utility in cohorts treated with next-generation AR signalling inhibitors (such as enzalutamide, apalutamide, or darolutamide) will determine whether these candidate biomarkers retain discriminatory power in modern therapeutic regimens following radical prostatectomy.
    
    \item \textbf{Exosomal circRNA liquid biopsy assays}: Given the marked stability of circRNAs in extracellular vesicles and bodily fluids \cite{Li2015Aug, Memczak2015}, profiling the circRNAs identified within our enzalutamide-resistance ceRNA network in patient plasma or urinary exosomes represents a minimally invasive strategy for longitudinal disease monitoring and early detection of biochemical relapse.

    \item \textbf{Therapeutic targeting of oncogenic circRNA-miRNA axes}: Unravelling specific circRNA hubs that sponge tumour-suppressive miRNAs in enzalutamide-resistant cells provides actionable targets for therapeutic intervention. Synthetic circRNA sponges or antisense oligonucleotides (ASOs) engineered to disrupt specific back-splice junctions could potentially re-sensitise resistant tumour cells to androgen receptor antagonism.
\end{itemize}

\subsection{Concluding Remarks}
This thesis embarked on a dual investigation into circular RNAs, guided by the central premise that purpose-built computational methods are essential to deconstruct non-coding RNA networks and uncover clinically relevant insights. The body of work presented herein has advanced this premise through a series of interconnected contributions that span from algorithmic benchmarking and software engineering to translational biomarker discovery in prostate cancer.

By deconstructing the post-transcriptional regulatory architecture of anti-androgen resistance, this thesis has demonstrated that circular RNAs are active participants within the competing endogenous RNA network. It has shown that transcriptional perturbations accompanying acquired enzalutamide resistance can be captured through coherent circRNA-miRNA-mRNA network axes, identifying regulatory interactions that are statistically informative for predicting biochemical recurrence across multi-cohort clinical datasets.

Crucially, this work has moved beyond descriptive characterisation to deliver actionable computational frameworks. The development of nf-core/circrna and pycircdb provides the scientific community with reproducible, containerised infrastructure that resolves long-standing obstacles in consensus quantification, pipeline portability, and cross-database annotation \cite{Digby2023Dec, Digby2024Dec}. Furthermore, the establishment of a 23-gene prognostic signature illustrates the utility of using regulatory RNA networks to constrain feature selection, providing a robust paradigm for biomarker prioritisation that captures a largely distinct prognostic space relative to existing clinical panels \cite{Mortensen2015Nov, Harrell1996Feb, Steyerberg2015Apr}.

In conclusion, this thesis underscores that circular RNAs are not artefactual splicing curiosities, but stable, functionally active regulators of gene expression with tangible relevance in oncology. The computational pipelines and network-modelling strategies delivered in this work establish a systematic framework for deciphering the role of non-coding RNAs. By demonstrating that non-coding regulatory interactions can be systematically modelled, harmonised across databases, and linked to patient outcomes in prostate cancer, this work provides a foundation for the continued integration of circular RNAs into precision medicine and translational oncology. 
